\chapter{Materials and Technologies}

This chapter introduces the core technologies used to implement the high-performance SDF computation pipeline. We describe the CUDA parallel computing platform, the NVIDIA OptiX ray tracing framework, and the dataset of 3D models used for evaluation.

\section{CUDA}

CUDA (Compute Unified Device Architecture) \cite{cuda} is NVIDIA's parallel computing platform. In this project, CUDA is used for kernel execution (launching massively parallel kernels for normal computation, SDF normalization, and anisotropic smoothing), memory management (allocating and managing GPU memory for mesh data, intermediate results, and output buffers), and the CUB library (using CUB's device-wide primitives for radix sort and unique operations when building the mesh adjacency graph in CSR format).

The key CUDA kernels implemented in this project include \texttt{GPUComputeBoundingBox} (computes the mesh bounding box using atomic min/max operations), \texttt{GPUComputeSDFMinMax} and \texttt{GPUApplySDFNormalization} (min-max normalization with logarithmic compression), \texttt{GPUGenerateEdges} and \texttt{GPUExtractCSR} (constructs the Compressed Sparse Row adjacency graph), and \texttt{AnisotropicSmoothingKernel} (bilateral smoothing kernel operating on the mesh graph).

\section{NVIDIA OptiX}

NVIDIA OptiX \cite{optix} is a high-level framework for ray tracing on NVIDIA GPUs. It provides hardware-accelerated ray tracing by leveraging NVIDIA RT Cores (available on Turing, Ampere, and newer architectures) to accelerate ray-triangle intersection tests in hardware. It automatically builds and maintains Bounding Volume Hierarchy (BVH) acceleration structures, enabling $O(\log F)$ ray-triangle intersection queries. The framework supports programmable shaders, allowing custom ray generation, miss, and closest-hit programs written in CUDA. It also enables zero-copy computation, where data remains on the GPU throughout the pipeline, eliminating host-device transfer overhead.

\begin{figure}[ht]
\centering
\begin{tikzpicture}[
    >=stealth,
    box/.style={
        rectangle,
        rounded corners=3pt,
        draw=red!60!black,
        fill=red!8,
        text width=5.5cm,
        minimum height=0.7cm,
        align=center,
        font=\small,
        inner sep=2pt
    },
    arr/.style={
        ->,
        -{stealth[length=2.5mm,width=2mm]}
    }
]

\node[box] (cuda) at (0, 0) {\textbf{Initialize CUDA}\\{\footnotesize Set device, query capabilities}};
\node[box] (ctx) at (0, -1.1) {\textbf{Create OptiX Context}\\{\footnotesize \texttt{optixDeviceContextCreate}}};
\node[box] (module) at (0, -2.2) {\textbf{Compile Module}\\{\footnotesize PTX $\to$ \texttt{optixModuleCreate}}};
\node[box] (prog) at (0, -3.3) {\textbf{Create Program Groups}\\{\footnotesize Raygen, Miss, Hitgroup}};
\node[box] (pipe) at (0, -4.4) {\textbf{Link Pipeline}\\{\footnotesize \texttt{optixPipelineCreate}}};
\node[box] (sbt) at (0, -5.5) {\textbf{Build Shader Binding Table}\\{\footnotesize Map shaders to geometry}};
\node[box] (bvh) at (0, -6.6) {\textbf{Build BVH}\\{\footnotesize \texttt{optixAccelBuild}}};
\node[box] (launch) at (0, -7.7) {\textbf{Launch Rays}\\{\footnotesize \texttt{optixLaunch}}};

\draw[arr] (cuda) -- (ctx);
\draw[arr] (ctx) -- (module);
\draw[arr] (module) -- (prog);
\draw[arr] (prog) -- (pipe);
\draw[arr] (pipe) -- (sbt);
\draw[arr] (sbt) -- (bvh);
\draw[arr] (bvh) -- (launch);

\end{tikzpicture}
\caption{Formal OptiX initialization sequence. Each step must be completed in order before rays can be traced.}
\label{fig:optix_init}
\end{figure}

The OptiX pipeline for SDF computation consists of module creation (compiling CUDA shader code to PTX and creating an OptiX module), program groups (defining ray generation, miss, and hit group programs), pipeline linking (connecting program groups into a traceable pipeline), Shader Binding Table (SBT) configuration (mapping shaders to geometry), BVH construction (building the acceleration structure from triangle mesh data), and launch (executing the ray tracing kernel across all vertices).

\section{Technology Selection}

The choice of OptiX + CUDA over alternative technologies was driven by several factors summarized in Table~\ref{tab:tech_compare}.

\begin{table}[h]
\centering
\begin{tabular}{@{}lll@{}}
\toprule
\textbf{Criterion} & \textbf{OptiX + CUDA} & \textbf{OpenGL (VCGlib)} \\
\midrule
Ray tracing & Hardware RT Cores & Software rasterization \\
BVH traversal & $O(\log F)$ hardware & None (depth peeling) \\
GPU requirement & NVIDIA RTX & Any GPU with OpenGL 3.3+ \\
Post-processing & Full GPU pipeline & None \\
Parallelism model & Per-vertex & Per-direction \\
\bottomrule
\end{tabular}
\caption{Comparison of OptiX vs.\ OpenGL approaches}
\label{tab:tech_compare}
\end{table}

As shown in Table~\ref{tab:tech_compare}, the OptiX + CUDA approach provides hardware-accelerated ray tracing through dedicated RT Cores, enabling $O(\log F)$ BVH traversal directly in hardware. The OpenGL approach relies on software rasterization with depth peeling and requires multiple camera directions to approximate thickness from all sides. OptiX's per-vertex parallelism maps naturally to the SDF computation (each vertex independently shoots rays), while OpenGL's per-direction parallelism requires $N$ separate draw calls with accumulated blending. The post-processing pipeline (log normalization and bilateral smoothing) runs entirely on the GPU in the OptiX approach, whereas the OpenGL approach has no built-in post-processing support.

\section{Dataset}

The evaluation dataset consists of 13 triangular mesh models in OBJ format, stored in the \texttt{Model/} directory. These models vary in complexity and size as shown in Table~\ref{tab:dataset}.

\begin{table}[h]
\centering
\begin{tabular}{@{}lrr@{}}
\toprule
\textbf{Model} & \textbf{Vertices} & \textbf{Description} \\
\midrule
360.obj & 2,200 & Simple object \\
9.obj & 2,639 & Simple object \\
400.obj & 3,703 & Medium complexity \\
76.obj & 5,923 & Medium complexity \\
181.obj & 7,242 & Medium complexity \\
118.obj & 9,153 & High complexity \\
368.obj & 11,202 & High complexity \\
112.obj & 13,628 & High complexity \\
369.obj & 13,606 & High complexity \\
158.obj & 14,587 & High complexity \\
371.obj & 14,599 & High complexity \\
Leaf.obj & 24,866 & Highest complexity \\
\bottomrule
\end{tabular}
\caption{Dataset models and vertex counts}
\label{tab:dataset}
\end{table}

As shown in Table~\ref{tab:dataset}, the dataset spans a wide range of mesh complexities, from simple objects like 360.obj with 2,200 vertices to highly detailed models like Leaf.obj with 24,866 vertices. This diversity ensures that the benchmark results are representative across different scales of geometric detail. The models are triangular meshes in OBJ format, which is the standard input format for the SDF computation pipeline. The vertex counts are used in Chapter 4 to analyze how computation time scales with model complexity.

The models were selected to cover a range of vertex counts (from 2,200 to 24,866) to evaluate how each implementation scales with mesh complexity.
